%------------------------------------
% DIPLOMARBEIT: Hauptdatei Dokument
% Version 1.2
%     
% Stefan Grünwald (2000 - 2007)
% Michael Fladischer (2012-2014)
%
% Das Dokument lässt sich am besten mit "latexmk" bauen:
% http://www.phys.psu.edu/~collins/software/latexmk-jcc/
%
% Die genutzten LaTeX-Packages sind alle Bestandteil der TeXlive Distribution die mit Debian (Squeeze) ausgeliefert wird.
% Für eine bequeme Installation unter Debian genügt es, folgende Zeile als "root" auszuführen:
%
% aptitude install -R latexmk texlive-base texlive-bibtex-extra texlive-binaries texlive-common texlive-fonts-extra texlive-fonts-recommended \
%  texlive-formats-extra texlive-generic-extra texlive-generic-recommended texlive-humanities texlive-lang-german texlive-latex-base texlive-latex-extra \
%  texlive-latex-recommended texlive-latex3 texlive-math-extra texlive-pictures texlive-plain-extra texlive-pstricks texlive-publishers texlive-science \
%  latex-beamer
%
% Alle .tex Dateien sind für den vim Editor ausgelegt und enthalten Modeline-Anweisungen zur besseren Darstellung.
% 
% Patches zur Verbesserung von Fehlern, zur Optimierung oder für neue Features können gerne an <michael@fladi.at> geschickt werden.
%------------------------------------

% Dokumentenklasse %{{{
\documentclass[
    %DIV12,                                 % Satzspiegelparameter 12
    %BCOR8mm,                               % Bindekorrektur 8mm
    %headinclude,                           % Berücksichtigung der Kopfzeile
                                            % bei der Ermittlung des Satzspiegels
    12pt,                                   % Schriftgrösse 12pt
    a4paper,                                % DIN A4
    captions=oneline,                       % 
    ngerman,
    draft=true,                             % Muss entfernt werden, bevor Arbeit abgegeben wird!
    bibliography=totoc,                     % Literaturverzeichniss automatisch im TOC einbinden
    listof=totoc,                           % Listen automatisch in TOC einbinden
    index=totoc]                            % Index automatisch in TOC einbinden
    {scrreprt}                              % KOMA-Script für Reports
%}}}

% Generische Packages %{{{
\usepackage[ngerman]{babel}
\usepackage[ngerman]{translator}
\usepackage[utf8]{inputenc}                 % Laden der UTF8-Eingabekodierung
                                            % ACHTUNG: Dateien müssen UTF8 codiert sein!
\usepackage[T1]{fontenc}                    % EC codierte Fonts
\usepackage{ae}                             % PDF encoding Info
\usepackage{palatino}                       % Brotschrift: Palationo
\usepackage{helvet}                         % Serifenlose Schrift: Helvetica
\usepackage[left]{eurosym}                  % EURO links
\usepackage{paralist}                       % Beeinflussung von Listen
\usepackage{tabularx}                       % Austreiben auf eine bestimmte Breite
\usepackage{longtable}                      % Mehrseitige Tabellen
\usepackage{ltxtable}                       % 'longtable' und 'tabularx'
\usepackage{array}                          % Zeilenumbruch in 'p'-Spalten
\usepackage{multirow}                       % Multiple Zeilen/Spalten in Tabellen
% http://www.ctan.org/tex-archive/macros/latex/contrib/booktabs/
% \toprule, \midrule und \bottomrule
\usepackage{booktabs}
% http://www.ctan.org/tex-archive/macros/latex/required/tools/
\usepackage{dcolumn}
% http://www.ctan.org/tex-archive/macros/latex/contrib/blindtext/
% Erzeugt Blindtext welcher beim Prüfen des Layouts behilflich sein kann
\usepackage{blindtext}
% http://ctan.org/tex-archive/macros/latex/contrib/siunitx/
%\usepackage{siunitx}
% http://www.ctan.org/tex-archive/macros/latex/required/graphics/
\usepackage{color}
% http://www.ctan.org/tex-archive/macros/latex/contrib/xcolor/
\usepackage[usenames,dvipsnames,svgnames,table]{xcolor}
% http://www.ctan.org/tex-archive/macros/latex/contrib/wasysym/
\usepackage[nointegrals]{wasysym}
% http://www.ctan.org/tex-archive/macros/latex/contrib/url/
\usepackage[hyphens]{url}
% http://www.ctan.org/tex-archive/macros/latex/contrib/mdwtools/
\usepackage{mdwlist}
% http://www.ctan.org/tex-archive/macros/latex/contrib/csquotes/
\usepackage{csquotes}
\usepackage{footnote}
\usepackage{footmisc}
\usepackage{gensymb}
\usepackage{subfigure}
\usepackage{wrapfig}
\usepackage{amsmath,amsthm,amssymb}
%\usepackage[geometry]{ifsym}
%\usepackage{scrhack}
%}}}

% Formatierung für Listings %{{{
% http://www.ctan.org/tex-archive/macros/latex/contrib/listings/
\usepackage{listings}
% Definition für JavaScript
\lstdefinelanguage{JavaScript}{
    comment=[l]{//},
    commentstyle=\color{green}\ttfamily,
    identifierstyle=\color{black},
    keywords={
        attributes,
        class,
        classend,
        do,
        empty,
        endif,
        endwhile,
        fail,
        function,
        functionend,
        if,
        implements,
        in,
        inherit,
        inout,
        not,
        of,
        operations,
        out,
        return,
        set,
        then,
        types,
        while,
        use,
    },
    keywordstyle=\color{blue}\bfseries,
    morecomment=[s]{/*}{*/},
    morestring=[b]",
    morestring=[b]',
    ndkeywords={}
    ndkeywordstyle=\color{yellow}\bfseries,
    sensitive=false,
    stringstyle=\color{red}\ttfamily,
}
\lstdefinelanguage{WKT}{
    keywordstyle=\color{blue}\bfseries,
    morekeywords={POINT, LINESTRING, POLYGON, MULTIPOINT, MULTILINESTRING, MULTIPOLYGON},
    morestring=[b](),
    sensitive=true,
}
\lstdefinelanguage{XML}
{
    identifierstyle=\color{blue},
    keywordstyle=\color{cyan},
    morecomment=[s]{<?}{?>},
    morekeywords={xmlns,version,type},
    morestring=[b]",
    morestring=[s]{>}{<},
    showstringspaces=false,
    stringstyle=\color{black},
}

% Sprachen laden
\lstloadlanguages{HTML}
% Format für Listings
\lstset{
    basicstyle=\ttfamily,
    breaklines=true,
    captionpos=b, % Beschriftungsposition am Ende des Listings
    columns=fixed,
    emphstyle=\bfseries\emph,
    escapeinside={@}{@},
    firstnumber=auto,
    float,
    frame=single, % einfacher Rahmen
    identifierstyle=\ttfamily,
    keywordstyle=,
    numbers=left,
    numbersep=5pt,
    numberstyle=\tiny,
    prebreak=\raisebox{0ex}[0ex][0ex]{\ensuremath{\hookleftarrow}},
    showspaces=false,
    showstringspaces=true,
    showtabs=false,
    stepnumber=2,
    xleftmargin=8pt,
}
%}}}

% Lokale Einstellungen %{{{
% vim: set tw=160:

\def\campusType{Studienarbeit}
\def\campusTitle{Software Engineering Methoden und ihre Auswahlkriterien je nach Entwicklungskontext}
\def\campusAuthor{Milena Berl}
\def\campusPersKnz{????}
\def\campusCity{Graz}
\def\campusDate{\today}
\def\campusSubject{}
\def\campusKeywords{PGP,x509,Firefox,XML-DSig}

%}}}

% Weiche für PDF/DVI %{{{
% http://www.ctan.org/tex-archive/macros/latex/contrib/oberdiek/
% Bitmaps stinken ja in der Druckstufe sowieso, besser TIKZ benutzen :-)
% für Encapsulated-Postscript und PDFTeX Grafiken

\usepackage[pdftex]{graphicx}
\graphicspath{{images/pdf/}}
\pdfcompresslevel=9
\usepackage{rotating}
%}}}

% TIKZ - Coole Grafiken direkt in TeX zeichnen :-) %{{{
% http://www.ctan.org/tex-archive/help/Catalogue/entries/pgf.html
\usepackage{tikz}
\usetikzlibrary{
    arrows,
    calc,
    decorations,
    decorations.pathreplacing,
    fadings,
    matrix,
    mindmap,
    patterns,
    positioning,
    shadows,
    shapes,
    trees,
}
% Globale Definitionen für TIKZ/PGF.
% Globale Definitionen für TIKZ/PGF.

% Styles %{{{
\tikzstyle{globe}=[ball color=white]
\tikzstyle{object}=[pattern=crosshatch dots,pattern color=black!25]
\tikzstyle{box}=[draw=red,dashed]
\tikzstyle{plane}=[draw=black,thick]
\tikzstyle{arrow}=[->,-latex',draw=red!75]
%}}}

% Functions %{{{
\newcommand\pgfmathsinandcos[3]{%
  \pgfmathsetmacro#1{sin(#3)}%
  \pgfmathsetmacro#2{cos(#3)}%
}
\newcommand\LongitudePlane[3][current plane]{%
  \pgfmathsinandcos\sinEl\cosEl{#2} % elevation
  \pgfmathsinandcos\sint\cost{#3} % azimuth
  \tikzset{#1/.estyle={cm={\cost,\sint*\sinEl,0,\cosEl,(0,0)}}}
}
\newcommand\LatitudePlane[3][current plane]{%
  \pgfmathsinandcos\sinEl\cosEl{#2} % elevation
  \pgfmathsinandcos\sint\cost{#3} % latitude
  \pgfmathsetmacro\yshift{\cosEl*\sint}
  \tikzset{#1/.estyle={cm={\cost,0,0,\cost*\sinEl,(0,\yshift)}}} %
}
\newcommand\DrawLongitudeCircle[2][1]{
  \LongitudePlane{\angEl}{#2}
  \tikzset{current plane/.prefix style={scale=#1}}
   % angle of "visibility"
  \pgfmathsetmacro\angVis{atan(sin(#2)*cos(\angEl)/sin(\angEl))} %
  \draw[current plane] (\angVis:1) arc (\angVis:\angVis+180:1);
  \draw[current plane,dashed] (\angVis-180:1) arc (\angVis-180:\angVis:1);
}
\newcommand\DrawLatitudeCircle[2][1]{
  \LatitudePlane{\angEl}{#2}
  \tikzset{current plane/.prefix style={scale=#1}}
  \pgfmathsetmacro\sinVis{sin(#2)/cos(#2)*sin(\angEl)/cos(\angEl)}
  % angle of "visibility"
  \pgfmathsetmacro\angVis{asin(min(1,max(\sinVis,-1)))}
  \draw[current plane] (\angVis:1) arc (\angVis:-\angVis-180:1);
  \draw[current plane,dashed] (180-\angVis:1) arc (180-\angVis:\angVis:1);
}
%}}}

% vim:filetype=tex foldmethod=marker autoindent expandtab shiftwidth=4 tw=160:
% Local variables:
% mode: latex
% End:              

%}}}

% Überschriften und Zeilenabstand %{{{
% http://www.ctan.org/tex-archive/macros/latex/required/tools/
\usepackage{rotating}
\usepackage{layout}

% http://www.ctan.org/tex-archive/macros/latex/contrib/caption/
% Überschriften formatieren
\usepackage[font=small,format=hang,labelfont=bf]{caption}

\renewcommand{\floatpagefraction}{.8}       % ,float page`, wenn diese mindestens
                                            % 80% der Seite einnehmen
                                            % default: .5
\renewcommand{\textfraction}{.1}            % neben den Abbildungen, die ,h`ere, ,t`op und
                                            % ,b`ottom plaziert werden, noch min. 10% Text
                                            % default: .2
\renewcommand{\topfraction}{.9}             % max. Anteil der Seite, die f\"{u}r Abbildungen,
                                            % die ,t`op plaziert werden, verwendet wird
                                            % default: .7
\renewcommand{\bottomfraction}{.6}          % max. Anteil der Seite, die f\"{u}r Abbildungen,
                                            % die ,b`ottom plaziert werden, verwendet wird
                                            % default: .3
%}}}

% Silbentrennung %{{{
% vim: set tw=160:

\hyphenation{Ge-schäfts-pro-zes-sen
             Ge-schäfts-pro-zess-mo-dell
             Ca-non-ic-aliz-a-tion-Meth-od
}


%}}}

% Index-Erstellung %{{{
\usepackage{makeidx}
\makeindex
%}}}

% PDF- und Hyperlink-Einstellungen %{{{
% http://www.ctan.org/tex-archive/macros/latex/contrib/hyperref/
\usepackage[plainpages=false,
            pdfpagelabels,
            pdftitle={\campusTitle},
            pdfauthor={\campusAuthor},
            pdfsubject={\campusType, \campusDate, \campusSubject},
            pdfkeywords={\campusKeywords},
            bookmarks=true,                 % Lesezeichen im 'pdf' analog zum Inhaltsverzeichnis,
            bookmarksopen=true,             % Bookmarks anzeigen...
            bookmarksnumbered=true,         % ...und nummerieren
            hyperindex=true,                % Indexeinträge als Hyperlinks
            colorlinks=false,               % Farbige Links
            linkcolor=black,
            urlcolor=black,
            citecolor=black,
            unicode=true,
            breaklinks,
            verbose
]{hyperref}
% http://www.ctan.org/tex-archive/macros/latex/contrib/breakurl/
% Allow links to have line-breaks even when processed by dvips/ps2pdf.
%\usepackage[hyphenbreaks]{breakurl}
%}}}

% APA style citations %{{{
% http://www.ctan.org/tex-archive/biblio/bibtex/contrib/apacite/
\usepackage{apacite}
%}}}

% Glossaries %{{{
% http://www.ctan.org/tex-archive/macros/latex/contrib/glossaries/
\usepackage[toc,acronym,style=long,nonumberlist]{glossaries} 
\makeglossaries
\loadglsentries{misc/glossary}
%}}}

% TODO-Notes %{{{
% http://www.ctan.org/tex-archive/macros/latex/contrib/todonotes/
\usepackage[colorinlistoftodos,obeyDraft]{todonotes} 
\reversemarginpar
%}}}

% Benutzerdefinierte Befehle %{{{
\newcommand{\xmlelem}[1]{\index{#1}\texttt{<#1/>}}
\newcommand{\xmlattr}[2]{\index{#1!#2}\textit{#2}}
\newcommand{\yes}{$\circledast$}
\newcommand{\no}{$\circleddash$}
\newcommand{\optional}{$\circledcirc$}
%\newcommand{\yes}{\rlap{\Cross}\Square}
%\newcommand{\no}{\Square}
%\newcommand{\optional}{\rlap{\FilledSmallCircle}\Square}
%\newcommand{\yes}{\rlap{\FilledSmallCircle}\Circle}
%\newcommand{\no}{\Circle}
%\newcommand{\optional}{\rlap{\SmallCircle}\Circle}
%\newcommand{\todo}{\marginpar{\textbf{TODO}}}
\renewcommand*{\glspostdescription}{}

%}}}

% Layout-Einstellungen %{{{

% Numerierungstiefe der Kapitel
\setcounter{secnumdepth}{3}%

% Numerierungstiefe des Inhaltsverzeichnisses
\setcounter{tocdepth}{3}%

% Abstand vor und nach Gleitobjekten im Text [h]
%\setlength{\intextsep}{12pt plus 2pt minus 2pt}    % default: {12pt plus 2pt minus 2pt}

% Abstand zwischen Gleitobjekten oben oder unten und Text [tb]
\setlength{\textfloatsep}{28pt plus 2pt minus 4pt}  % default: {20pt plus 2pt minus 4pt}


%%%---- PARAGRAPH ---------------------
\makeatletter
\renewcommand{\paragraph}{%
  \@startsection {paragraph}{4}{\z@}%
  {-3.25ex plus-1ex minus-.2ex}{1sp}%
  {\raggedsection \normalfont \sectfont \nobreak \size@paragraph \nobreak } % was: \reset@font \normalsize\bf
} \makeatother

%%%---- SUBPARAGRAPH ------------------
\makeatletter
\renewcommand{\subparagraph}{%
  \@startsection {subparagraph}{5}{\parindent}%
  {3.25ex \@plus 1ex \@minus .2ex}{-1em}%
  {\raggedsection \normalfont \sectfont \nobreak \normalsize\bf \nobreak } % was: \size@subparagraph
} \makeatother

\clubpenalty = 10000
\widowpenalty = 10000

%%%------------------------------------
%%%---- SEITENDEKLARATION -------------
%%%------------------------------------
\setlength{\textheight}{660pt}%
\setlength{\voffset}{40pt}%
\setlength{\hoffset}{20pt}%
%\setlength{\footskip}{20pt}%
%}}}

% Kopf- und Fußzeilen %{{{
% http://www.ctan.org/tex-archive/macros/latex/contrib/fancyhdr/
\usepackage{fancyhdr}
\pagestyle{fancy}
% Redefine fetching of chapter name
\renewcommand{\chaptermark}[1]{\markboth{\thechapter.\ #1}{}}
% Page header
\lhead{}
\chead{}
\rhead{\small\leftmark}
% Page footer
\lfoot{}
\cfoot{\small\thepage}
\rfoot{}
%}}}

% Smarte Referenzen %{{{
% http://www.ctan.org/tex-archive/macros/latex/contrib/fancyref/
\usepackage[german]{cleveref}
\usepackage{nameref}
%}}}

% Haupt-Dokument %{{{

\begin{document}

% Titelseite
\begin{titlepage}
%\begin{sf}
\begin{sffamily}
\begin{center}
\huge \campusType \\

\vspace{2cm}%

\campusTitle \\
\vspace{3cm}%

\begin{normalsize}

ausgeführt am

\begin{figure} [h!]
    \centerline{\includegraphics[width=0.5\textwidth]{Campus02_Logo}}
\end{figure}

%\vspace{4.5cm}%

Studiengang\\ 
Informationstechnologien und Wirtschaftsinformatik\\

\vspace{3.5cm}%

%\begin{tabularx}{\linewidth}{ll}
Von: \campusAuthor \\
Pers. Kennz. \campusPersKnz
%\end{tabularx}

\vspace{3.5cm}%
\end{normalsize}
\end{center}

\begin{tabularx}{0.9\linewidth}{Xc}
\campusCity, am \campusDate & ............................. \\
         & \campusAuthor\\
\end{tabularx}

\end{sffamily}
\end{titlepage}

% Ende der Titelseite
\newpage

% Liste der TODOS (nur sichtbar wenn die TODOs auch aktiviert sind, siehe Package todonotes)
\listoftodos

% Seitennummerierung auf Römisch setzen und bei Seite 1 beginnen
\pagenumbering{roman}
\setcounter{page}{1}

\setlength{\parindent}{0em}

% Ehrenwort
% vim: set tw=160:

\chapter*{Ehrenwörtliche Erklärung}

Ich erkläre ehrenwörtlich, dass ich die vorliegende Arbeit selbstständig und ohne fremde Hilfe verfasst, andere als die angegebenen Quellen nicht benützt und
die benutzten Quellen wörtlich zitiert sowie inhaltlich entnommene Stellen als solche kenntlich gemacht habe.\\

\vspace{1 cm} \hfill
\begin{tabular*}{4.4 cm}{c}
.................................... \\
\campusAuthor \\
\end{tabular*}


% Danksagung
% vim: set tw=160:

\chapter*{Danksagung}

\vspace{1.5cm}

\begin{center}
\raggedright
Ich danke meinem größten Schatz Markus Bohlayer, für seine stete Geduld und Unterstützung. In der Hitze Siziliens hat er sich meinen Fragen, meinem Frust und meiner Angst vor italienischen Kitesurfern mit Bravour gestellt. Ohne Ihn wäre mir LaTex ein Mysterium geblieben. 
\ \\
\raggedleft{\campusAuthor}
\end{center}

\vspace{1.5cm}

\campusCity, am \campusDate


% Kurzfassung / Abstract
% vim: set tw=160:

\chapter*{Kurzfassung}
\raggedright
Im Rahmen dieser Studienarbeit werden die verbreitesten Methoden der Softwareentwicklung, so wie in der Literatur vorzufinden, dargestellt. Ziel ist es, die Faktoren, welche zur Auswahl der richtigen Methode je nach Entwicklungskontext, darzustellen. 

% Hack damit zwei Chapter auf eine Seite passen
\begingroup
\let\clearpage\relax
\chapter*{Abstract}
\endgroup
\blindtext


% Inhaltsverzeichnis
\tableofcontents

% Hauptteil
\setlength{\parskip}{3ex plus0.5ex minus0.5ex}

% Seitennummerierung auf Arabisch setzen und bei Seite 1 beginnen
\clearpage
\pagenumbering{arabic}
\setcounter{page}{1}

% Laden des eigentlichen Inhalts
% vim: set tw=160:

% vim: set tw=160:

\chapter{Einleitung}
\label{chap:Einleitung}
\blindtext \todo{Ein\-leit\-ung verfassen}

\section{Ausgangssituation}
\label{sec:Einleitung:ausgangssituation}
\blindtext

\section{Aufgabenstellung}
\label{sec:Einleitung:aufgabenstellung}
Die primäre Aufgabenstellung dieser Arbeit besteht darin, \ldots

Folgende Aspekte sollen deshalb in dieser Arbeit erläutert werden:
\begin{itemize}
    \item Wie?
    \item Wo?
    \item Wann?
\end{itemize}

\section{Zielsetzung der Arbeit}
\label{sec:Einleitung:ziele}
\blindtext \todo{Ziele definieren}

\section{Vorgehen und Methodik}
\label{sec:Einleitung:methodik}
\blindtext

\section{Aufbau der Arbeit}
\label{sec:Einleitung:aufbau}
\blindtext

% vim: set tw=160:

\chapter{Grundlagen und Definitionen}
\label{chap:GrundlagenDefinitionen}
\blindtext

Eine Referenz auf die \cref{fig:wot} oder auf die \cref{fig:ht} ist nicht schwer und nimmt, falls nötig, sogar die Seitenzahl mit ein den Text auf. Wer Mut und
Laune hat, kann sich auch mit TIKZ\footnote{\url{http://www.texample.net/tikz/}} beschäftigen um schönere Grafiken als mit Bitmaps zu erhalten.

\begin{figure}
\centering
\begin{tikzpicture}[node distance   = 4 cm]
    \tikzset{Key/.style = {shape          = rectangle,
                           rounded corners,
                           drop shadow,
                           fill           = white,
                           draw           = black,
                           thick,
                           text           = black,
                           text centered,
                           inner sep      = 3pt,
                           outer sep      = 0pt,
                           minimum size   = 24 pt}}
    \tikzset{DirectTrust/.style   = {thick,
                                     double          = black,
                                     double distance = 1pt}}
    \tikzset{IndirectTrust/.style   = {thin,
                                       double          = white,
                                       double distance = 1pt}}
    \tikzset{Label/.style =   {draw,
                               fill           = white,
                               text           = black}}
    \node[Key](A){Alice};
    \node[Key,right=of A](B){Bob};
    \node[Key,right=of B](C){Carmen};
    \node[Key,above=of B](D){Dorian};     
    \draw[DirectTrust](B) to node[Label]{vertraut} (D) ;
    \tikzset{DirectTrust/.append style = {bend left}}
    \draw[DirectTrust](A) to node[Label]{vertraut} (B);
    \draw[DirectTrust](B) to node[Label]{vertraut} (C);
    \draw[DirectTrust](D) to node[Label]{vertraut} (C);
    \tikzset{IndirectTrust/.append style = {bend left}}
    \draw[IndirectTrust](A) to node[Label]{vertraut indirekt} (D) ;
    \tikzset{IndirectTrust/.append style = {bend right}}
    \draw[IndirectTrust](A) to node[Label]{vertraut indirekt} (C) ;

\end{tikzpicture}

\caption{Schematische Darstellung des \glsentrydesc{wot}}
\label{fig:wot}
\end{figure}


% Beispiel für Rotation einer Abbildung, für weitere Optionen siehe:
% https://en.wikibooks.org/wiki/LaTeX/Rotations
\begin{figure}
\centering
\begin{sideways}
\begin{tikzpicture}[node distance   = 3 cm]
    \tikzset{Key/.style = {shape          = rectangle,
                           rounded corners,
                           drop shadow,
                           fill           = white,
                           draw           = black,
                           thick,
                           text           = black,
                           text centered,
                           inner sep      = 3pt,
                           outer sep      = 0pt,
                           minimum size   = 24 pt}}
    \tikzset{DirectTrust/.style   = {thick,
                                     double          = black,
                                     double distance = 1pt}}
    \tikzset{IndirectTrust/.style   = {thin,
                                       double          = white,
                                       double distance = 1pt}}
    \tikzset{Label/.style =   {draw,
                               fill           = white,
                               text           = black,
                               text centered,
                               text width     = 4em}}
    \node[Key](CA1){CA1 (Root)};
    \node[below=of CA1](AUX1){};
    \node[Key,left=of AUX1](CA1-1){CA1.1};
    \node[Key,right=of AUX1](CA1-2){CA1.2};
    \node[below=of CA1-1](AUX1-1){};
    \node[Key,left=of AUX1-1](Alice){Alice};
    \node[Key,right=of AUX1-1](Bob){Bob};
    \node[Key,below=of CA1-1](Carmen){Carmen};
    \node[Key,below=of CA1-2](Dorian){Dorian};
    \draw[DirectTrust,bend left](CA1-1) to node[Label]{vertraut} (CA1) ;
    \draw[DirectTrust,bend right](CA1-2) to node[Label]{vertraut} (CA1) ;
    \draw[DirectTrust,bend left](Alice) to node[Label]{vertraut} (CA1-1) ;
    \draw[DirectTrust,bend right](Bob) to node[Label]{vertraut} (CA1-1) ;
    \draw[DirectTrust](Carmen) to node[Label]{vertraut} (CA1-1) ;
    \draw[DirectTrust](Dorian) to node[Label]{vertraut} (CA1-2) ;
    \draw[IndirectTrust,bend right](Carmen) to node[Label]{vertraut indirekt} (Dorian) ;
    \draw[IndirectTrust,bend right](Alice) to node[Label]{vertraut indirekt} (Carmen) ;

\end{tikzpicture}

\end{sideways}
\caption{Schematische Darstellung eines \glsentrydesc{ht}}
\label{fig:ht}
\end{figure}

Es kann auch aus dem \cref{listing:link-element} auf die Zeile \ref{line:link-href} verwiesen werden, ohne dies fix im Text zu inkludieren. Alles wird automatisch über die
Position der Labels berechnet.

% Source-Code Listings aus externer Datei, siehe auch:
% https://en.wikibooks.org/wiki/LaTeX/Source_Code_Listings
\lstinputlisting[language=HTML,
                 caption={Erklärender Text},
                 label=listing:link-element,
                 emph={link,href}]{source/link.html}

Auch Tabellen wie die \cref{tab:xml-namespaces} sind schnell erstellt und sehr gut lesbar. Vor allem sollte auf vertikale Trennlinien verzichtet werden, um die
Lesbarkeit zu erhöhen \cite{latex}. Jetzt kommt \gls{xml}. Wie \citeA{niag} gezeigt hat \ldots

% Tabelle mit tabularx-Umgebung für definierte Breite, siehe:
% https://en.wikibooks.org/wiki/LaTeX/Tables
\begin{table}
    \begin{center}
        \begin{tabularx}{.8\textwidth}{ l X }
            \gls{xml}-Dialekt  & \gls{xml-ns} \\
            \hline
            \glstext{xhtml} & \url{http://www.w3.org/1999/xhtml} \\
            \glstext{xml-dsig} & \url{http://www.w3.org/2000/09/xmldsig\#} \\
            \glstext{xinclude} & \url{http://www.w3.org/2001/XInclude} \\
            \glstext{xslt} & \url{http://www.w3.org/1999/XSL/Transform} \\
            \hline
        \end{tabularx}
    \end{center}
    \caption{Namespaces für unterschiedliche \protect\gls{xml}-Dialekte}
    \label{tab:xml-namespaces}
\end{table}

Platzhalter für Abbildungen sind auch leicht möglich und integrieren sich gut mit der TODO-Liste am Anfang des Dokuments. Als Beispiel siehe
\cref{fig:missing}.

\begin{figure}
    \begin{center}
        \missingfigure[figwidth=.7\textwidth]{Schöne Abbildung mit TIKZ zeichnen.}
    \end{center}
    \caption{Ein Platzhalter für eine Abbildung}
    \label{fig:missing}
\end{figure}


% vim: set tw=160:

\chapter{Zusammenfassung}
\label{chap:Zusammenfassung}
Um nun ein schönes PDF-Dokument zu erhalten muss nur noch folgender Befehl im Verzeichnis mit \texttt{bachelorarbeit.tex} und \texttt{presentation.tex}
ausgeführt werden: \todo[color=green!40]{PDF kompilieren}

\begin{verbatim}
latexmk
\end{verbatim}

Danach befinden sich \texttt{bachelorarbeit.pdf} sowie \texttt{presentation.pdf} im Verzeichnis. Um nach dem Bauen der Dokumente das Verzeichnis wieder
aufzuräumen genügt folgender Befehl:

\begin{verbatim}
latexmk -c
\end{verbatim}

% vim: set tw=160:

\chapter{Grundlagen und Definitionen}
\label{chap:GrundlagenDefinitionen}
\blindtext

Eine Referenz auf die \cref{fig:wot} oder auf die \cref{fig:ht} ist nicht schwer und nimmt, falls nötig, sogar die Seitenzahl mit ein den Text auf. Wer Mut und
Laune hat, kann sich auch mit TIKZ\footnote{\url{http://www.texample.net/tikz/}} beschäftigen um schönere Grafiken als mit Bitmaps zu erhalten.

\begin{figure}
\centering
\begin{tikzpicture}[node distance   = 4 cm]
    \tikzset{Key/.style = {shape          = rectangle,
                           rounded corners,
                           drop shadow,
                           fill           = white,
                           draw           = black,
                           thick,
                           text           = black,
                           text centered,
                           inner sep      = 3pt,
                           outer sep      = 0pt,
                           minimum size   = 24 pt}}
    \tikzset{DirectTrust/.style   = {thick,
                                     double          = black,
                                     double distance = 1pt}}
    \tikzset{IndirectTrust/.style   = {thin,
                                       double          = white,
                                       double distance = 1pt}}
    \tikzset{Label/.style =   {draw,
                               fill           = white,
                               text           = black}}
    \node[Key](A){Alice};
    \node[Key,right=of A](B){Bob};
    \node[Key,right=of B](C){Carmen};
    \node[Key,above=of B](D){Dorian};     
    \draw[DirectTrust](B) to node[Label]{vertraut} (D) ;
    \tikzset{DirectTrust/.append style = {bend left}}
    \draw[DirectTrust](A) to node[Label]{vertraut} (B);
    \draw[DirectTrust](B) to node[Label]{vertraut} (C);
    \draw[DirectTrust](D) to node[Label]{vertraut} (C);
    \tikzset{IndirectTrust/.append style = {bend left}}
    \draw[IndirectTrust](A) to node[Label]{vertraut indirekt} (D) ;
    \tikzset{IndirectTrust/.append style = {bend right}}
    \draw[IndirectTrust](A) to node[Label]{vertraut indirekt} (C) ;

\end{tikzpicture}

\caption{Schematische Darstellung des \glsentrydesc{wot}}
\label{fig:wot}
\end{figure}


% Beispiel für Rotation einer Abbildung, für weitere Optionen siehe:
% https://en.wikibooks.org/wiki/LaTeX/Rotations
\begin{figure}
\centering
\begin{sideways}
\begin{tikzpicture}[node distance   = 3 cm]
    \tikzset{Key/.style = {shape          = rectangle,
                           rounded corners,
                           drop shadow,
                           fill           = white,
                           draw           = black,
                           thick,
                           text           = black,
                           text centered,
                           inner sep      = 3pt,
                           outer sep      = 0pt,
                           minimum size   = 24 pt}}
    \tikzset{DirectTrust/.style   = {thick,
                                     double          = black,
                                     double distance = 1pt}}
    \tikzset{IndirectTrust/.style   = {thin,
                                       double          = white,
                                       double distance = 1pt}}
    \tikzset{Label/.style =   {draw,
                               fill           = white,
                               text           = black,
                               text centered,
                               text width     = 4em}}
    \node[Key](CA1){CA1 (Root)};
    \node[below=of CA1](AUX1){};
    \node[Key,left=of AUX1](CA1-1){CA1.1};
    \node[Key,right=of AUX1](CA1-2){CA1.2};
    \node[below=of CA1-1](AUX1-1){};
    \node[Key,left=of AUX1-1](Alice){Alice};
    \node[Key,right=of AUX1-1](Bob){Bob};
    \node[Key,below=of CA1-1](Carmen){Carmen};
    \node[Key,below=of CA1-2](Dorian){Dorian};
    \draw[DirectTrust,bend left](CA1-1) to node[Label]{vertraut} (CA1) ;
    \draw[DirectTrust,bend right](CA1-2) to node[Label]{vertraut} (CA1) ;
    \draw[DirectTrust,bend left](Alice) to node[Label]{vertraut} (CA1-1) ;
    \draw[DirectTrust,bend right](Bob) to node[Label]{vertraut} (CA1-1) ;
    \draw[DirectTrust](Carmen) to node[Label]{vertraut} (CA1-1) ;
    \draw[DirectTrust](Dorian) to node[Label]{vertraut} (CA1-2) ;
    \draw[IndirectTrust,bend right](Carmen) to node[Label]{vertraut indirekt} (Dorian) ;
    \draw[IndirectTrust,bend right](Alice) to node[Label]{vertraut indirekt} (Carmen) ;

\end{tikzpicture}

\end{sideways}
\caption{Schematische Darstellung eines \glsentrydesc{ht}}
\label{fig:ht}
\end{figure}

Es kann auch aus dem \cref{listing:link-element} auf die Zeile \ref{line:link-href} verwiesen werden, ohne dies fix im Text zu inkludieren. Alles wird automatisch über die
Position der Labels berechnet.

% Source-Code Listings aus externer Datei, siehe auch:
% https://en.wikibooks.org/wiki/LaTeX/Source_Code_Listings
\lstinputlisting[language=HTML,
                 caption={Erklärender Text},
                 label=listing:link-element,
                 emph={link,href}]{source/link.html}

Auch Tabellen wie die \cref{tab:xml-namespaces} sind schnell erstellt und sehr gut lesbar. Vor allem sollte auf vertikale Trennlinien verzichtet werden, um die
Lesbarkeit zu erhöhen \cite{latex}. Jetzt kommt \gls{xml}. Wie \citeA{niag} gezeigt hat \ldots

% Tabelle mit tabularx-Umgebung für definierte Breite, siehe:
% https://en.wikibooks.org/wiki/LaTeX/Tables
\begin{table}
    \begin{center}
        \begin{tabularx}{.8\textwidth}{ l X }
            \gls{xml}-Dialekt  & \gls{xml-ns} \\
            \hline
            \glstext{xhtml} & \url{http://www.w3.org/1999/xhtml} \\
            \glstext{xml-dsig} & \url{http://www.w3.org/2000/09/xmldsig\#} \\
            \glstext{xinclude} & \url{http://www.w3.org/2001/XInclude} \\
            \glstext{xslt} & \url{http://www.w3.org/1999/XSL/Transform} \\
            \hline
        \end{tabularx}
    \end{center}
    \caption{Namespaces für unterschiedliche \protect\gls{xml}-Dialekte}
    \label{tab:xml-namespaces}
\end{table}

Platzhalter für Abbildungen sind auch leicht möglich und integrieren sich gut mit der TODO-Liste am Anfang des Dokuments. Als Beispiel siehe
\cref{fig:missing}.

\begin{figure}
    \begin{center}
        \missingfigure[figwidth=.7\textwidth]{Schöne Abbildung mit TIKZ zeichnen.}
    \end{center}
    \caption{Ein Platzhalter für eine Abbildung}
    \label{fig:missing}
\end{figure}



% Anhang

% Abkürzungen
\clearpage
\printglossaries

% Abbildungen
\clearpage
\listoffigures

% Tabellen
\clearpage
\listoftables

% Listings
\clearpage
\lstlistoflistings

% Literaturverzeichnis
\clearpage
\bibliographystyle{apacite}
\bibliography{misc/literatur}

% Index
\clearpage
\printindex

% Das tiefgründige Zitat, optional :-)
% vim: set tw=160:

\chapter*{}


\mbox{}
\renewcommand{\dictumwidth}{\textwidth}
\dictum[Roy Trenneman, The IT Crowd]{Have you tried turning it off and on again?}
%

\end{document}
%}}}

% vim:filetype=tex foldmethod=marker autoindent expandtab shiftwidth=4 tw=160:
% Local variables:
% mode: latex
% End:              
